\section{Special Rules}

\subsection{Reserve}

The Necrons strike suddenly and vanish without leaving a trace. Their mastery of teleportation gives them tactical opportunities for redeployment and an optimal distribution of their forces.

During army deployment, any detachment may be placed in Reserve so long as it has a means of arriving on the battlefield, such as Portals large enough for it or special rules or abilities such as Deep Strike. Detachments in Reserve receive orders normally, but may activate only to join the field during their activation.

A detachment in Reserve that no longer has a means of entering the battlefield is considered destroyed.

\subsection{Phase Out}

Necron commanders follow their own patterns and sometimes vanish without warning. When a Necron detachment with several bases is Broken, it must make a Phase Out roll. To do so, roll 1D6: on a 2+ nothing happens; otherwise they disappear from the battlefield and are considered destroyed. If there is a character attached to the detachment, this roll automatically succeeds.

\subsection{Engines of Destruction}

Only one Engine of Destruction is allowed in a Necron army per 5000 points. The Engines of Destruction are the Abattoir and the Aeonic Orb.

\subsection{Scarabs}

Canoptek Scarabs are small beetle-like robots, constantly repairing damaged Necrons. Moving in a silent hover through the empty corridors of the Tomb World's necropolis, they are among the countless mechanical servants used for maintenance for millennia.

Scarab tokens may be used by any Necron base of class 1 or 2. Each token may be used in two ways.

\textbf{Assault:} A class 1 or 2 base directs the Scarabs that attack the enemy. In addition to its ranged weapons, the base gains a one-use shot that uses a Template (7.5~cm), hits on 5+, has a range of 20~cm, no AP, and the Reduces Cover ($-1$) ability. A base may benefit from only one Scarab token per turn in this way.

\textbf{Repair:} when a base makes a Reanimation Protocols roll, you may use a Scarab token to improve the base's chance of success by 2.
